\chapter{引言}

随着物联网（Internet of Things, IoT）技术的快速发展，越来越多的智能设备接入网络，物联网安全问题日益突出。物联网设备通常具有资源受限、部署环境复杂、数量庞大等特点，使得传统的安全技术难以直接应用于物联网场景。因此，研究适用于物联网环境的安全技术具有重要的理论意义和实际价值。

本文旨在综述对称加密技术在物联网领域的最新进展。对称加密作为一种高效的加密方式，在资源受限的物联网设备中具有广泛的应用前景。通过对近年来相关文献的调研和分析，本文将从轻量级对称加密算法、物联网安全协议、以及实际应用案例等方面进行系统性的总结和讨论。

\section{研究背景}

物联网是指通过各种信息传感设备，按照约定的协议，将任何物品与互联网连接起来，进行信息交换和通信，以实现智能化识别、定位、跟踪、监控和管理的一种网络。物联网的应用场景涵盖了智能家居、智慧城市、工业控制、医疗健康等多个领域。

然而，物联网的快速发展也带来了严峻的安全挑战。物联网设备通常部署在无人值守的环境中，容易受到物理攻击；设备的计算能力和存储资源有限，难以运行复杂的加密算法；此外，物联网设备数量庞大，密钥管理也成为一大难题。

\section{研究意义}

对称加密技术是物联网安全的基础技术之一。与公钥加密相比，对称加密具有计算效率高、实现简单等优势，特别适合资源受限的物联网设备。近年来，随着物联网应用的不断扩展，针对物联网环境的轻量级对称加密算法研究取得了显著进展。

本报告通过对称加密技术在物联网领域最新进展的综述，有助于研究人员和工程技术人员全面了解该领域的研究现状和发展趋势，为物联网安全系统的设计和实现提供参考。
